% Compile in this directory with: latexmk -pdf SVP_consistency_sharp.tex
\documentclass[pdflatex,sn-mathphys-num]{sn-jnl}
\usepackage{amsmath,amssymb,amsfonts,amsthm,booktabs,array}
\DeclareMathOperator*{\argmin}{arg\,min}
\newcommand{\Prob}{\mathbb P}
\newcommand{\E}{\mathbb E}
\newcommand{\R}{\mathbb R}
\newcommand{\Z}{\mathbb Z}
\newcommand{\C}{\mathcal C}
\newcommand{\taus}{\tau^\star}
\newcommand{\Kh}{\widehat K_n}
\newcommand{\tht}{\widehat\tau_n}
\newcommand{\tol}{\widetilde\tau_n}
\newcommand{\op}{o_{\Prob}}
\newcommand{\Op}{O_{\Prob}}
\newcommand{\norm}[1]{\left\lVert#1\right\rVert}
\newcommand{\abs}[1]{\left\lvert#1\right\rvert}
\theoremstyle{thmstyleone}
\newtheorem{theorem}{Theorem}
\newtheorem{proposition}[theorem]{Proposition}
\newtheorem{lemma}[theorem]{Lemma}
\newtheorem{corollary}[theorem]{Corollary}
\theoremstyle{thmstyletwo}
\newtheorem{remark}{Remark}
\raggedbottom
\begin{document}
\title[Consistency and multiscale localization of SVP]{Sharp consistency, multiscale localization and oracle asymptotics for Gaussian smallest valid partitioning}
\author*[1]{\fnm{Vincent} \sur{Runge}}\email{vincent.runge@univ-evry.fr}
\affil*[1]{\orgdiv{Laboratoire de Math\'ematiques et Mod\'elisation d'Evry},
\orgname{Universit\'e Paris-Saclay, CNRS, Univ Evry},
\city{Evry-Courcouronnes}, \postcode{91037}, \country{France}}
\abstract{We study Gaussian smallest valid partitioning with a least-squares tie-break, for both the ordinary single-change scan and a multiscale subsegment constraint. For the ordinary scan, a threshold just above the second-order Darling--Erd\H{o}s boundary yields consistent selection on macroscopic segments; its leading scale is $\log\log n$. A central plateau demonstrates failure on shorter segments despite strong local signals. Testing all subsegments with a correction based on the effective length of each split removes this dilution. A simultaneous CUSUM bound gives uniform finite-sample overestimation control, and local witnesses give explicit count and localization bounds for unequal lengths, shrinking jumps and increasing segment counts. The corrected null maximum is bounded in probability, so an arbitrarily slowly diverging additive threshold suffices under the corresponding signal condition. For fixed nonzero jumps on macroscopic segments, both procedures equal the least-squares estimator supplied with the true count with probability tending to one; their location, level and empirical-loss limits follow. We distinguish improved detection conditions from the irreducible integer location uncertainty and specify the validity and pruning requirements of the multiscale construction.}
\keywords{Change-point detection, multiscale inference, Gaussian bridge, consistency, oracle equivalence, nonregular asymptotics}
\maketitle

\section{Setting and principal result}\label{sec:setting}

For each $n$, let
\begin{equation}\label{eq:model}
Y_i=\mu_i+\sigma\varepsilon_i,\qquad
\varepsilon_1,\ldots,\varepsilon_n\ \text{i.i.d. }\mathcal N(0,1),
\qquad \sigma>0.
\end{equation}
The means are deterministic and constant on the segments of
$\taus=(0=\taus_0<\taus_1<\cdots<\taus_{K_n^\star}=n)$.
Write $m_{k,n}=\taus_k-\taus_{k-1}$, $\ell_n=\min_k m_{k,n}$,
and $\theta_{k,n}$ for the mean on segment $k$. All internal changes are genuine:
$\Delta_{j,n}=|\theta_{j+1,n}-\theta_{j,n}|>0$.
When $K_n^\star\ge2$, put $\Delta_n=\min_j\Delta_{j,n}$.
The no-change case $K_n^\star=1$ is treated separately whenever this minimum occurs.
No dependence assumption between different rows $n$ is needed for convergence in probability.

For integer boundaries $0\le a<b\le n$, $Y_{a..b}$ comprises observations
$a+1,\ldots,b$. With known $\sigma$, define
\begin{align}
\C(Y_{a..b})&=\frac{1}{2\sigma^2}\sum_{i=a+1}^b(Y_i-\overline Y_{a..b})^2,\label{eq:cost}\\
Z_u(a,b)&=\frac{1}{\sigma}\sqrt{\frac{(u-a)(b-u)}{b-a}}
  (\overline Y_{a..u}-\overline Y_{u..b}),\label{eq:contrast}\\
f(Y_{a..b})&=\frac12\max_{a<u<b} Z_u(a,b)^2.\label{eq:scan}
\end{align}
Set $f=0$ for a singleton. A partition $\tau$ is $\gamma$-valid if
every segment has $f\le\gamma$. The SVP estimator $\tht$ minimizes
$(K(\tau),Q(\tau))$ in lexicographic order over valid partitions, where
$Q(\tau)=\sum_k\C(Y_{\tau_{k-1}..\tau_k})$ and $K$ counts \emph{segments}.
Singletons ensure feasibility for every $\gamma\ge0$.
Use lexicographically smallest internal boundary vectors to resolve any remaining ties.
Sections~\ref{sec:setting}--\ref{sec:limits} concern the exact optimizer with
\texttt{subtests = "none"}. Section~\ref{sec:multiscale} defines a separate
multiscale estimator with \texttt{subtests = "both"}; its scale correction
is part of the validity statistic.

For all sufficiently large $n$, define
\begin{equation}\label{eq:centering}
d_n=\log\log n+\frac12\log\log\log n-\frac12\log\pi.
\end{equation}
Values for finitely many small $n$ may be assigned arbitrarily.

\begin{theorem}[Consistency at the sharp null threshold]\label{thm:main}
Suppose $K_n^\star=K^\star$ is fixed and $\ell_n\ge an$ for a fixed $a>0$.
Let $\gamma_n\ge0$ be deterministic and satisfy
\begin{equation}\label{eq:threshold-main}
\gamma_n-d_n\longrightarrow+\infty.
\end{equation}
If $K^\star\ge2$, assume also
\begin{equation}\label{eq:signal-main}
\gamma_n=o(n\Delta_n^2/\sigma^2).
\end{equation}
Then
\begin{equation}\label{eq:count-main}
\Prob(\Kh=K^\star)\longrightarrow1.
\end{equation}
For $K^\star\ge2$, on the event of equal counts match boundaries in increasing order.
For every $\eta>0$,
\begin{equation}\label{eq:relative-main}
\Prob\left(\Kh=K^\star,\ 
\max_{1\le j<K^\star}|\widehat\tau_j-\taus_j|<\eta n\right)
\longrightarrow1.
\end{equation}
Under $K^\star=1$, condition~\eqref{eq:threshold-main} is also necessary for
\eqref{eq:count-main}.
\end{theorem}

Thus $\gamma_n=(1+\epsilon)\log\log n$, for any fixed $\epsilon>0$,
suffices whenever $n\Delta_n^2/(\sigma^2\log\log n)\to\infty$.
More finely, one may use $\gamma_n=d_n+w_n$ with $w_n\to\infty$ arbitrarily slowly,
provided~\eqref{eq:signal-main} holds. The necessity statement concerns uniform
calibration over a class including constant signals; it does not assert necessity
of every signal-strength condition below. All thresholds use the cost units
in~\eqref{eq:cost}.

The proof is given in Section~\ref{sec:consequences}. Its main ingredients are a
finite-sample Gaussian bridge bound and a deterministic implication of validity.
The latter also covers increasing segment counts and unequal segment lengths.
For the multiscale estimator, Theorems~\ref{thm:ms-finite} and~\ref{thm:ms-consistency}
give recovery guarantees under local signal conditions, while Theorem~\ref{thm:ms-oracle}
preserves the fixed-signal oracle limits. Throughout, $K$ counts segments.
Section~\ref{sec:oracle} strengthens the conclusion to oracle equivalence and
nonregular location limits when the segment levels are fixed.

\section{Null calibration: finite samples and the critical boundary}\label{sec:null}

For $m\ge2$ and $q>1$, set
\begin{equation}\label{eq:peeling-number}
\begin{aligned}
J_q(m)&=\max\left\{1,\left\lceil\frac{2\log(m-1)}{\log q}\right\rceil\right\},\\
U_m(\gamma)&=\min\left\{1,(m-1)e^{-\gamma},
\inf_{q>1}2J_q(m)e^{-\gamma/q}\right\}.
\end{aligned}
\end{equation}
Set $U_1(\gamma)=0$. The infimum is a deterministic bound, not a data-dependent
choice of threshold.

\begin{lemma}[Gaussian bridge bound]\label{lem:null}
On a constant-mean segment of length $m$, for every $\gamma\ge0$,
\begin{equation}\label{eq:null-tail}
\Prob\{f(Y_{a..a+m})>\gamma\}\le U_m(\gamma).
\end{equation}
For each fixed $q>1$ there is $C_q<\infty$ such that
$U_m(\gamma)\le C_q(1+\log m)e^{-\gamma/q}$ for every $m\ge2$.
\end{lemma}
\begin{proof}
Let $S_k=\sum_{i=1}^k\varepsilon_{a+i}$. The standardized contrasts are
$(S_k-kS_m/m)/\sqrt{k(1-k/m)}$. Their joint distribution is exactly
$B(k/m)/\sqrt{(k/m)(1-k/m)}$, with $B$ a standard Brownian bridge.
The covariance identity
$B(t)\overset{d}=(1-t)W(t/(1-t))$ gives the process representation
\[
\left\{\frac{B(t)}{\sqrt{t(1-t)}}:0<t<1\right\}
\overset{d}=\left\{\frac{W(s)}{\sqrt{s}}:s=\frac{t}{1-t}>0\right\},
\]
where $W$ is standard Brownian motion. The discrete scan is bounded by the
continuous supremum over $s\in[(m-1)^{-1},m-1]$.
Cover this interval with $J_q(m)$ intervals $[v,qv]$, starting at
$v=(m-1)^{-1}$. On such an interval,
\[
\Prob\left\{\sup_{v\le s\le qv}\frac{|W(s)|}{\sqrt{s}}>\sqrt{2\gamma}\right\}
\le\Prob\left\{\sup_{0\le s\le qv}|W(s)|>\sqrt{2\gamma v}\right\}
\le4\overline\Phi(\sqrt{2\gamma/q})\le2e^{-\gamma/q}.
\]
The reflection principle gives the penultimate inequality.
The last one uses $2\overline\Phi(x)\le e^{-x^2/2}$, $x\ge0$.
Union over the covering intervals proves the peeling bound.
Alternatively, the $m-1$ discrete contrasts are each standard normal, giving
$(m-1)e^{-\gamma}$. Taking the minimum proves the claim, including $m=2$.
\end{proof}

Let $p_m(\gamma)$ denote the exact null probability
$\Prob\{f(Y_{0..m})\le\gamma\}$. It depends on $m$ and $\gamma$ only.
Independence between true segments gives
\begin{equation}\label{eq:true-feasibility}
\Prob\{\taus\text{ is valid}\}=\prod_{k=1}^{K_n^\star}p_{m_{k,n}}(\gamma),
\qquad
\Prob(\Kh>K_n^\star)\le1-\prod_kp_{m_{k,n}}(\gamma)
\le\sum_k U_{m_{k,n}}(\gamma).
\end{equation}
No independence between contrasts within one segment is assumed.

\begin{proposition}[Second-order null boundary]\label{prop:DE}
For a constant-mean segment of length $m\to\infty$,
\begin{equation}\label{eq:DE}
f(Y_{0..m})-d_m\ \Longrightarrow\ H,
\qquad \Prob(H\le x)=\exp\{-2e^{-x}\},\quad x\in\R.
\end{equation}
Consequently, for arbitrary deterministic $\gamma_n\ge0$ in the no-change model,
\begin{equation}\label{eq:null-iff}
\Prob(\Kh=1)\to1\quad\Longleftrightarrow\quad\gamma_n-d_n\to+\infty.
\end{equation}
If $\gamma_n-d_n\to x\in\R$, this probability tends to $\exp\{-2e^{-x}\}$;
if $\gamma_n-d_n\to-\infty$, it tends to zero.
\end{proposition}
\begin{proof}
Put $M_m=\max_{1\le k<m}|S_k-kS_m/m|/\sqrt{k(1-k/m)}$.
The scalar Darling--Erd\H{o}s theorem states that
\[
\sqrt{2\log\log m}\,M_m-
\left(2\log\log m+\tfrac12\log\log\log m-\tfrac12\log\pi\right)
\Longrightarrow H.
\]
An accessible theorem-level source is \cite[Equation (3.7), p.~8, and
Theorem 3.10, p.~11]{torgovitski2016}: take dimension one and independent
standard Gaussian observations (equivalently, constant functions with that
random value). The covariance eigenvalue is one; stationarity, every required
moment, causality, and the approximation condition hold, the latter with zero
error. This verifies the hypotheses rather than using their displayed
Assumption (L) without justification.
Squaring the expansion gives $M_m^2/2=d_m+H_m+\op(1)$ with
$H_m\Longrightarrow H$, because $(\log\log\log m)^2/\log\log m\to0$.
This proves~\eqref{eq:DE}.

There is exactly one one-segment partition, so $\{\Kh=1\}=
\{f(Y_{0..n})\le\gamma_n\}$. Tightness proves sufficiency in~\eqref{eq:null-iff}.
If $\gamma_n-d_n$ does not tend to $+\infty$, along a subsequence it is bounded
above by some $A<\infty$; the limsup along that subsequence is at most
$\exp\{-2e^{-A}\}<1$. This proves necessity. The remaining assertions follow
from continuity and the tails of the limit distribution.
\end{proof}

\begin{remark}[Fixed error probability]
The asymptotic null threshold at rejection probability $\alpha\in(0,1)$ is
$d_n+\log2-\log\{-\log(1-\alpha)\}$. It has a limiting nonzero false-split
probability and therefore does not give count consistency. For fixed
$K^\star$ and $m_{k,n}/n\to p_k>0$, the true partition is valid at threshold
$d_n+x$ with limiting probability $\exp\{-2K^\star e^{-x}\}$.
This last probability is a feasibility probability, not a claimed exact
selection probability when $K^\star>1$.
\end{remark}

\section{A deterministic argument and finite-sample recovery}\label{sec:finite}

\begin{lemma}[What validity forces locally]\label{lem:local}
Let $(a,b]$ be $\gamma$-valid, of length $M=b-a$. If $t,r$ are integers,
$r\ge1$, and $a\le t-r<t<t+r\le b$, then
\begin{equation}\label{eq:local-constraint}
\frac{1}{\sigma}\sqrt{\frac r2}
  |\overline Y_{t..t+r}-\overline Y_{t-r..t}|
\le2\sqrt{\frac{\gamma M}{r}}.
\end{equation}
\end{lemma}
\begin{proof}
Define $V(u)=\sum_{i=a+1}^u(Y_i-\overline Y_{a..b})$, including $V(a)=V(b)=0$.
Validity and~\eqref{eq:contrast} imply, at every integer $u\in[a,b]$,
\[
|V(u)|\le\sigma\sqrt{2\gamma}\sqrt{\frac{(u-a)(b-u)}{M}}
\le\sigma\sqrt{\frac{\gamma M}{2}}.
\]
Since
$r(\overline Y_{t..t+r}-\overline Y_{t-r..t})=V(t+r)-2V(t)+V(t-r)$,
the triangle inequality and division by $\sigma\sqrt{2r}$ give the result.
The argument remains valid when $t-r=a$ or $t+r=b$.
\end{proof}

Suppress $n$ in the next theorem. For each true change choose a deterministic
integer radius $r_j\ge1$ with
\begin{equation}\label{eq:radii}
r_j\le\min(m_j,m_{j+1}),\qquad
r_j+r_{j+1}\le m_{j+1}+1\quad(1\le j<K^\star-1).
\end{equation}
The integer neighborhoods $\mathcal B_j=\{t\in\Z:|t-\taus_j|<r_j\}$
are then disjoint and lie in $\{1,\ldots,n-1\}$.
The uniform choice $r_j=\lceil\ell_n/2\rceil$ is always allowed, even when
$\ell_n=1$. Define
\begin{equation}\label{eq:margin}
A_j=\frac{\Delta_j}{\sigma}\sqrt{\frac{r_j}{2}},\qquad
B_j=2\sqrt{\frac{\gamma n}{r_j}},\qquad
v_j=\Phi(B_j-A_j)-\Phi(-B_j-A_j).
\end{equation}
Here $\Phi$ is the standard normal distribution function.

\begin{theorem}[Finite-sample count and location bounds]\label{thm:finite}
Under~\eqref{eq:model}, for every $\gamma\ge0$ and radii~\eqref{eq:radii},
\begin{align}
\Prob(\Kh>K^\star)&\le1-\prod_kp_{m_k}(\gamma),\label{eq:over}\\
\Prob(\Kh<K^\star)&\le\sum_{j=1}^{K^\star-1}v_j,\label{eq:under}\\
\Prob\left(\Kh=K^\star,\ |\widehat\tau_j-\taus_j|<r_j\ \forall j\right)
&\ge\prod_kp_{m_k}(\gamma)-\sum_jv_j.\label{eq:joint}
\end{align}
In particular, the failure probability in~\eqref{eq:joint} is bounded by
\begin{equation}\label{eq:failure-bound}
\sum_k U_{m_k}(\gamma)+
\sum_j\exp\left\{-\frac12(A_j-B_j)_+^2\right\},
\qquad x_+=\max(x,0).
\end{equation}
Empty sums and products have their usual meanings, so the result includes
$K^\star=1$. Bounds may always be truncated to $[0,1]$.
\end{theorem}
\begin{proof}
On the deterministic windows around $\taus_j$, let
\[
D_j=\frac{1}{\sigma}\sqrt{\frac{r_j}{2}}
  (\overline Y_{\taus_j..\taus_j+r_j}
    -\overline Y_{\taus_j-r_j..\taus_j}).
\]
The windows lie inside the two adjacent true segments. Thus $D_j$ is normal
with variance one and mean of absolute value $A_j$, and
$\Prob(|D_j|\le B_j)=v_j$.
Suppose $|D_j|>B_j$. Any partition without an internal boundary in
$\mathcal B_j$ contains $[\taus_j-r_j,\taus_j+r_j]$ in the closure of a single
segment. That segment has length at most $n$ and is invalid by
Lemma~\ref{lem:local}. Therefore, on $\bigcap_j\{|D_j|>B_j\}$, every valid
partition has a boundary in each $\mathcal B_j$, and hence at least $K^\star$
segments. If it has exactly $K^\star$ segments, there is exactly one boundary
in each neighborhood, with the asserted ordered matching.

When the true partition is valid, SVP has at most $K^\star$ segments.
Intersect these events and use the union bound for the $D_j$. Their mutual
independence is unnecessary. This proves~\eqref{eq:over}--\eqref{eq:joint}.
Finally, if $A_j>B_j$, symmetry and a one-sided normal tail give
$v_j\le\Phi(B_j-A_j)\le\exp\{-(A_j-B_j)^2/2\}$.
If $A_j\le B_j$, use $v_j\le1$; combine with Lemma~\ref{lem:null}.
\end{proof}

\begin{corollary}[Explicit confidence radii]\label{cor:radii}
Let $K^\star\ge2$, $\alpha\in(0,1)$, and choose admissible radii satisfying
\begin{equation}\label{eq:confidence-radii}
r_j\ge\max\left\{
\frac{4\sqrt{2}\,\sigma\sqrt{n\gamma}}{\Delta_j},
\frac{16\sigma^2}{\Delta_j^2}\log\frac{K^\star-1}{\alpha}
\right\}.
\end{equation}
The probability in~\eqref{eq:joint} is at least
$1-\sum_kU_{m_k}(\gamma)-\alpha$.
\end{corollary}
\begin{proof}
The first inequality implies $B_j\le A_j/2$; the second implies
$A_j^2/8=\Delta_j^2r_j/(16\sigma^2)\ge\log((K^\star-1)/\alpha)$.
Apply~\eqref{eq:failure-bound}.
\end{proof}

These radii are guarantees from validity alone. They are not the final
statistical precision of the least-squares tie-break. Their dependence on
$\sqrt n$ reflects the ability of a whole-interval scan to dilute a local
change. Section~\ref{sec:limits} shows that dilution can be a genuine
obstruction to selection for plain SVP.

\section{Consistency consequences}\label{sec:consequences}

\begin{proof}[Proof of Theorem~\ref{thm:main}]
Since $an\le m_{k,n}\le n$, we have $d_{m_{k,n}}-d_n\to0$ uniformly over the
fixed number of true segments. Proposition~\ref{prop:DE} and
$\gamma_n-d_n\to\infty$ imply true-partition feasibility with probability
tending to one. This argument applies to every sequence of segment lengths
in $[an,n]$: failure of uniform tightness would supply a contradicting
subsequence to~\eqref{eq:DE}.

When $K^\star\ge2$, fix $0<b<\min(a/2,\eta)$ and take $r_j=\lfloor bn\rfloor$.
These radii satisfy~\eqref{eq:radii} for large $n$.
Uniformly in $j$,
\[
\frac{B_j}{A_j}=\frac{2\sqrt{2}\sigma\sqrt{n\gamma_n}}{\Delta_{j,n}r_j}
\longrightarrow0,
\qquad
\min_j A_j^2\ge\frac{\Delta_n^2\lfloor bn\rfloor}{2\sigma^2}\longrightarrow\infty.
\]
The latter follows because $\gamma_n\to\infty$ and~\eqref{eq:signal-main} holds.
Theorem~\ref{thm:finite} proves both count and relative-location consistency.
For $K^\star=1$, use Proposition~\ref{prop:DE} directly.
\end{proof}

\begin{corollary}[Increasing segment counts]\label{cor:growing}
Allow $K_n^\star\ge2$ to depend on $n$. Suppose $\gamma_n\to\infty$ and
\begin{equation}\label{eq:growing}
\sum_k U_{m_{k,n}}(\gamma_n)\to0,\qquad
\frac{\ell_n^2\Delta_n^2}{\sigma^2 n\gamma_n}\to\infty,\qquad
\log(K_n^\star+1)=o(\ell_n\Delta_n^2/\sigma^2).
\end{equation}
Then $\Prob(\Kh=K_n^\star)\to1$.
If also $\ell_n\to\infty$, then, for every $\eta>0$,
\[
\Prob\left(\Kh=K_n^\star,\
\max_j|\widehat\tau_j-\taus_j|<\eta\ell_n\right)\to1.
\]
\end{corollary}
\begin{proof}
For count recovery use $r_j=\lceil\ell_n/2\rceil$. Then $B_j/A_j\to0$ uniformly
and, eventually, the second term in~\eqref{eq:failure-bound} is at most
$(K_n^\star-1)\exp\{-\ell_n\Delta_n^2/(32\sigma^2)\}$, which tends to zero.
For relative localization use $r_j=\lfloor b\ell_n\rfloor$ with
$0<b<\min(1/2,\eta)$. The same argument applies with a different fixed constant.
\end{proof}

The first condition in~\eqref{eq:growing} is explicit and accommodates different
calibrations. For example, it follows from
$K_n^\star(1+\log n)e^{-\gamma_n/q}\to0$ for a fixed $q>1$.
If a deterministic bound $K_n^\star\le\overline K_n$ is available, one choice is
$\gamma_n=q\{\log\overline K_n+\log(1+\log n)+w_n\}$ with $w_n\to\infty$.
Without such a bound, $\gamma_n=c\log n$, $c>1$, gives
$\sum_kU_{m_{k,n}}(\gamma_n)\le n^{1-c}$.
The remaining conditions in~\eqref{eq:growing} must still be checked.
They are sufficient conditions for this estimator, not minimax conditions
for arbitrary change-point procedures.

\begin{proposition}[Almost-sure count recovery]\label{prop:strong}
Place the rows on any common probability space. Suppose $K^\star$ is fixed,
$\ell_n\ge an$, and, when $K^\star\ge2$, $\inf_n\Delta_n/\sigma>0$.
For $\gamma_n=c\log n$ with any $c>1$, almost surely
$\Kh=K^\star$ for every sufficiently large $n$.
Conversely, for a single nested i.i.d. Gaussian sequence with constant mean,
if $\gamma_n\le c\log n$ eventually for some $0\le c<1$, then
$\Kh>1$ infinitely often almost surely.
\end{proposition}
\begin{proof}
Choose $1<q<c$. The overestimation bound is
$O((1+\log n)n^{-c/q})$, which is summable. With radii proportional to $n$,
the underestimation bound in~\eqref{eq:failure-bound} is $O(e^{-c_1n})$ for
some $c_1>0$. The first Borel--Cantelli lemma gives the first assertion.
There is no assumption of independence between rows.

For the converse, standardize to $\sigma=1$. The split at $n-1$ gives
$f(Y_{0..n})\ge\tfrac12(n-1)n^{-1}
(\varepsilon_n-\overline\varepsilon_{1:n-1})^2$.
Choose $c<c'<1$. The independent events
$\varepsilon_n^2>2c'\log n$ occur infinitely often, since their probabilities
have divergent sum. The strong law gives
$\overline\varepsilon_{1:n-1}\to0$; along these events the displayed gain
exceeds $c\log n$ eventually.
\end{proof}

This distinguishes convergence in probability at the $\log\log n$ scale from
eventual almost-sure correctness at the $\log n$ scale. The borderline
coefficient $c=1$ is not classified by Proposition~\ref{prop:strong}.

\section{Oracle equivalence and sharp statistical consequences}\label{sec:oracle}

In this section impose the classical fixed-signal regime:
\begin{equation}\label{eq:classical}
K_n^\star=K^\star<\infty,\qquad
\taus_j/n\to t_j,\quad 0=t_0<t_1<\cdots<t_{K^\star}=1,
\qquad \theta_{k,n}=\theta_k,
\end{equation}
where $\theta_{k+1}\ne\theta_k$ and $\sigma>0$ is fixed.
The oracle partition $\tol$ minimizes $Q(\tau)$ over all partitions with
exactly $K^\star$ segments. ``Oracle'' refers only to knowing the count;
its boundary locations and segment means are estimated from the data.

\begin{lemma}[Oracle localization]\label{lem:oracle-tight}
Under~\eqref{eq:classical},
$\max_j|\widetilde\tau_j-\taus_j|=\Op(1)$.
For $K^\star=1$ this is interpreted as an empty maximum.
\end{lemma}
The full argument, including the local likelihood expansion and its
global localization input, is in Appendix~\ref{app:oracle}.

\begin{lemma}[Feasibility after bounded endpoint errors]\label{lem:edited}
Under~\eqref{eq:classical}, fix an integer $M\ge0$. For every deterministic
interval with endpoints within $M$ of the endpoints of one true segment,
\begin{equation}\label{eq:edited}
f(Y_{a..b})-d_n\quad\text{is bounded above in probability},
\end{equation}
uniformly over this finite family of intervals.
Consequently, if $\gamma_n-d_n\to\infty$, the oracle partition is valid
with probability tending to one.
\end{lemma}
\begin{proof}
Each interval has length proportional to $n$. Subtract its associated true
segment level. Its deterministic mean vector is nonzero at at most $2M$
observations, with uniformly bounded magnitudes. Write each contrast as
$Z_u=N_u+L_u$, where $N_u$ is the pure-noise contrast on that interval and
$L_u$ is its mean. The normalized contrast vector has Euclidean norm one,
so $\sup_u|L_u|\le C_M$. If both sides of the split have length at least $h$,
its coefficients have absolute values at most $h^{-1/2}$, giving
$|L_u|\le C_M/\sqrt h$.

Take $h_n=\lceil(\log\log n)^4\rceil$.
For splits with both lengths at least $h_n$, Proposition~\ref{prop:DE} gives
$\max_u|N_u|=\Op(\sqrt{\log\log n})$, hence uniformly
$\tfrac12(N_u+L_u)^2\le\tfrac12N_u^2+\op(1)$.
The at most $2h_n$ remaining splits have
$\max|N_u|=\Op(\sqrt{\log(h_n+1)})$ by a Gaussian union bound.
Their observed squared contrasts are $\Op(\log(h_n+1)+1)=\op(\log\log n)$.
The pure-noise maximum on the whole interval equals
$d_{b-a}+\Op(1)=d_n+\Op(1)$, proving~\eqref{eq:edited}.
There are at most $K^\star(2M+1)^2$ intervals, so the bound is uniform.

Lemma~\ref{lem:oracle-tight} puts all oracle endpoints in this family with
probability arbitrarily close to one by choosing $M$ large, independently
of $n$. For fixed $M$,~\eqref{eq:edited} and $\gamma_n-d_n\to\infty$ give
validity with probability tending to one. Let $M\to\infty$ after $n\to\infty$.
\end{proof}

\begin{theorem}[Equality with the oracle estimator]\label{thm:oracle}
Under~\eqref{eq:classical}, if
\begin{equation}\label{eq:oracle-threshold}
\gamma_n-d_n\to\infty,\qquad \gamma_n=o(n),
\end{equation}
then
\begin{equation}\label{eq:oracle-equality}
\Prob(\tht=\tol)\longrightarrow1.
\end{equation}
In particular, SVP has exactly the same asymptotic location and level
distributions as least squares supplied with the true segment count.
\end{theorem}
\begin{proof}
Theorem~\ref{thm:main} gives count recovery. Lemma~\ref{lem:edited} gives
oracle feasibility. On their intersection, SVP minimizes $Q$ on a subset
of the $K^\star$-segment partitions that contains the unrestricted minimizer.
It therefore equals that minimizer, with the common tie convention.
In fact the oracle minimizer is almost surely unique: for distinct
$K$-segment partitions the orthogonal projections onto their piecewise
constant subspaces are distinct. Equality of their residual sums of squares
is the zero set of a nonzero quadratic polynomial in a Gaussian vector
having a density. A finite union of these zero sets has probability zero.
\end{proof}

For explicit limits, put $\delta_j=(\theta_{j+1}-\theta_j)/\sigma$, including
its sign. Let $\{\xi_{j,r}^{+},\xi_{j,r}^{-}:j=1,\ldots,K^\star-1,r\ge1\}$
be independent standard normal variables. Define independent two-sided
random walks, with $R_j(0)=0$, by
\begin{equation}\label{eq:random-walk}
R_j(h)=
\begin{cases}
\displaystyle\sum_{r=1}^{h}(\delta_j^2/2+\delta_j\xi_{j,r}^{+}),&h>0,\\[0.3em]
\displaystyle\sum_{r=1}^{-h}(\delta_j^2/2-\delta_j\xi_{j,r}^{-}),&h<0.
\end{cases}
\qquad H_j=\argmin_{h\in\Z}R_j(h).
\end{equation}
Each walk tends to $+\infty$ in both directions almost surely, and its
minimum is almost surely unique. Denote fitted SVP segment means by
$\widehat\theta_k$ and its fitted mean vector by $\widehat\mu$.
For limits involving indexed boundaries and levels, assign arbitrary
values on $\{\Kh\ne K^\star\}$; this convention does not affect convergence
in distribution or in probability.

\begin{theorem}[Joint limits and empirical loss]\label{thm:limits}
Under the assumptions of Theorem~\ref{thm:oracle},
\begin{equation}\label{eq:joint-limit}
\left((\widehat\tau_j-\taus_j)_{j<K^\star},
\left(\frac{\sqrt{m_{k,n}}}{\sigma}(\widehat\theta_k-\theta_k)\right)_{k\le K^\star}\right)
\Longrightarrow ((H_j)_{j<K^\star},(G_k)_{k\le K^\star}),
\end{equation}
where the $G_k$ are independent standard normal variables, independent of
all the $H_j$. Moreover,
\begin{equation}\label{eq:loss-limit}
\frac{1}{\sigma^2}\sum_{i=1}^n(\widehat\mu_i-\mu_i)^2
\Longrightarrow
\sum_{k=1}^{K^\star}G_k^2+
\sum_{j=1}^{K^\star-1}\delta_j^2|H_j|.
\end{equation}
Consequently, $\max_j|\widehat\tau_j-\taus_j|=\Op(1)$ and
$n^{-1}\norm{\widehat\mu-\mu}^2=\Op(\sigma^2/n)$ for fixed $K^\star$.
\end{theorem}
\begin{proof}
The oracle location limit and its joint independence from the normalized
true-segment noise sums are proved in Appendix~\ref{app:oracle}.
Bounded-in-probability endpoint errors alter each segment sum by $\Op(1)$
and its length by $\Op(1)$. Since $m_{k,n}\asymp n$,
\begin{equation}\label{eq:means-expansion}
\widetilde\theta_k-\theta_k
=\frac{\sigma}{m_{k,n}}\sum_{i=\taus_{k-1}+1}^{\taus_k}\varepsilon_i
+\Op(n^{-1}).
\end{equation}
For completeness, the $\Op(1)$ assertion is obtained by first restricting
endpoint errors to a fixed $M$, taking the maximum over the finitely many
involved boundary observations, and then letting $M\to\infty$.
The normalized sums on disjoint true segments are exactly independent
standard normals. This gives the level limit jointly with the location limit.

On correctly assigned observations of true segment $k$, the sum of squared
errors is $(m_{k,n}+\Op(1))(\widetilde\theta_k-\theta_k)^2$.
At boundary $j$, exactly $|\widetilde\tau_j-\taus_j|$ observations are
assigned to the neighboring level. Their contribution is
$(\theta_{j+1}-\theta_j)^2|\widetilde\tau_j-\taus_j|+\op(1)$.
The boundary neighborhoods are disjoint with probability tending to one.
Summing proves the oracle version of~\eqref{eq:loss-limit}; apply
Theorem~\ref{thm:oracle} to transfer all conclusions to SVP.
These are distributional and stochastic-order statements; no convergence
of expectations is inferred without uniform integrability.
\end{proof}

\begin{remark}[An information bound for exact localization]
Even if both levels and the segment count are known, uniform exact
localization is impossible at a fixed jump. Consider two models with a
single boundary at $t$ and $t+1$, respectively, with the same levels on
either side and standardized jump $\delta\ne0$. Their data laws differ in
one Gaussian observation only. The minimal sum of the two testing errors
is $2\Phi(-|\delta|/2)$, obtained by comparing that observation with the
midpoint of the two levels. Any boundary estimator induces a test (choose
the first model exactly when the estimate is $t$), and therefore
\[
\max\{\Prob_t(\widehat\tau\ne t),
\Prob_{t+1}(\widehat\tau\ne t+1)\}\ge\Phi(-|\delta|/2)>0.
\]
Thus the nonvanishing integer location uncertainty in~\eqref{eq:joint-limit}
reflects a statistical obstruction even in this simpler experiment.
\end{remark}

\begin{remark}[Exact recovery and F1]
For a fixed nonzero jump, $\Prob(R_j(1)<0)=\Phi(-|\delta_j|/2)>0$, so
$\Prob(H_j\ne0)>0$. Thus count consistency does not imply exact integer
recovery. For a fixed integer matching tolerance $d\ge0$ and $K^\star\ge2$,
the limiting expected F1 is
\[
\lim_{n\to\infty}\E[\mathrm{F1}_n(d)]
=\frac{1}{K^\star-1}\sum_{j=1}^{K^\star-1}\Prob(|H_j|\le d).
\]
Indeed, with probability tending to one, the count is correct and the
macroscopically separated boundaries can only match their corresponding
true boundaries. Integer-valued convergence gives the displayed matching
probabilities, and boundedness permits passage to expectations.
If instead $d_n'\to\infty$, F1 tends to one in probability and in expectation.
The notation $d_n'$ here is unrelated to the centering~\eqref{eq:centering}.
\end{remark}

\section{Relation with penalized least squares}\label{sec:coupling}

\begin{proposition}[Exact comparison at a common threshold]\label{prop:coupling}
Let $\widehat\tau^{\mathrm{OP}}$ minimize $Q(\tau)+\beta K(\tau)$, where
$\beta=\gamma\ge0$. Every OP segment is $\gamma$-valid, and
$\Kh\le\widehat K^{\mathrm{OP}}$. If the counts agree, the partitions agree
under the common tie convention. Hence, for every statistic $T$ taking
values in $[0,1]$,
\begin{align}\label{eq:coupling}
\E|T(\tht)-T(\widehat\tau^{\mathrm{OP}})|
&\le\Prob(\tht\ne\widehat\tau^{\mathrm{OP}})\nonumber\\
&\le\Prob(\Kh\ne K^\star)+\Prob(\widehat K^{\mathrm{OP}}\ne K^\star).
\end{align}
This proposition is deterministic apart from its final probability statement
and holds for any additive cost whose split gains define the validity scan.
\end{proposition}
\begin{proof}
If an OP segment had a split gain larger than $\beta$, splitting it would
decrease $Q+\beta K$, a contradiction. OP is therefore feasible for SVP,
which proves the count inequality. At equal counts, minimizing $Q$ first
over all partitions and then over the valid subset containing the OP
minimizer gives equality of costs and, with common tie resolution, of
partitions. The probability bound and bounded-statistic comparison follow.
\end{proof}

At a sufficiently large common logarithmic threshold, known penalized
least-squares consistency and Theorem~\ref{thm:main} imply asymptotic equality
of the entire partitions. This is stronger than equality of F1 scores.
In contrast, the sharp $\log\log n$ calibration of SVP does not require
penalized least squares to be consistent at that same penalty.

If OP uses residual sum of squares without the factor $1/(2\sigma^2)$,
its matching penalty is $\beta_{\mathrm{RSS}}=2\sigma^2\gamma$.
A result requiring a sufficiently large coefficient of $\log n$ does not
by itself validate a particular BIC coefficient. Neither this coupling nor
the null calibration in Section~\ref{sec:null} automatically applies to
\texttt{right} or \texttt{both} subsegment tests.

\section{Limits of the result and a variance extension}\label{sec:limits}

\begin{proposition}[Cancellation on a central plateau]\label{prop:plateau}
Let $n,m_n$ have the same parity, $0<m_n<n$, and
$L_n=(n-m_n)/2$. Consider the three-segment mean vector
\[
\mu_i=\Delta\,\mathbf1\{L_n<i\le L_n+m_n\},\qquad \Delta\ne0.
\]
Its noiseless whole-series scan is exactly
\begin{equation}\label{eq:plateau}
f(\mu_{0..n})=
\frac{\Delta^2m_n^2}{2\sigma^2n}\frac{n-m_n}{n+m_n}.
\end{equation}
If $\gamma_n\ge(1+\epsilon)\log\log n$ eventually for some $\epsilon>0$ and
$\Delta^2m_n^2/(\sigma^2n)=o(\gamma_n)$, then
$\Prob(\Kh=1)\to1$, although $K_n^\star=3$.
Furthermore, if
$\Delta^2m_n^2\log\log n/(\sigma^2n)\to0$, then
$f(Y_{0..n})-d_n\Longrightarrow H$ with the same law as in~\eqref{eq:DE}.
In this latter regime, every deterministic threshold satisfying
$\gamma_n-d_n\to\infty$ yields $\Prob(\Kh=1)\to1$.
\end{proposition}
\begin{proof}
Let $V_\mu(u)=\sum_{i=1}^u(\mu_i-\Delta m_n/n)$.
Before the plateau $V_\mu(u)=-u\Delta m_n/n$, on the plateau
$V_\mu(u)=\Delta(1-m_n/n)(u-n/2)$, and after it
$V_\mu(u)=\Delta m_n(1-u/n)$.
The squared standardized contrast is $nV_\mu(u)^2/(\sigma^2u(n-u))$.
It increases up to $u=L_n$ before the plateau, is maximal at one of the
two plateau endpoints within the plateau, and decreases after the plateau.
Evaluation at either endpoint proves~\eqref{eq:plateau}.

The triangle inequality for the maximum absolute contrast gives
$\sqrt{f(Y_{0..n})}\le\sqrt{f(\sigma\varepsilon_{0..n})}
+\sqrt{f(\mu_{0..n})}$.
Proposition~\ref{prop:DE} implies
$f(\sigma\varepsilon_{0..n})/\log\log n\to1$ in probability.
The assumptions leave a fixed multiplicative gap between the noise scan
and $\gamma_n$, while $f(\mu)=o(\gamma_n)$.
Thus the whole series is valid with probability tending to one, and SVP
then selects its one-segment partition.
For the last assertion, the reverse triangle inequality also gives
\[
|f(Y_{0..n})-f(\sigma\varepsilon_{0..n})|
\le2\sqrt{f(\sigma\varepsilon_{0..n})f(\mu_{0..n})}+f(\mu_{0..n})=\op(1).
\]
Apply Proposition~\ref{prop:DE} and tightness.
\end{proof}

For example, take even $n$ and $m_n=2\lfloor n^{1/4}/2\rfloor$, with fixed
$\Delta/\sigma\ne0$. Then
$m_n\Delta^2/(\sigma^2\log n)\to\infty$, but plain SVP selects one segment
at every deterministic threshold that is consistent under a constant mean.
This is an analytic counterexample to consistency under only a local
spacing--jump condition. A smaller threshold sufficient for null
consistency does not eliminate this cancellation phenomenon.

\begin{corollary}[A controlled plug-in variance]\label{cor:variance}
Use a common positive estimate $\widehat\sigma_n^2$ in both cost and validity.
Independence from the data is not required. The conclusions of
Theorem~\ref{thm:main} remain true if its deterministic signal condition holds,
$\widehat\sigma_n^2/\sigma^2\to1$ in probability, and
\begin{equation}\label{eq:variance-margin}
\gamma_n\left|\widehat\sigma_n^2/\sigma^2-1\right|
=\op(\gamma_n-d_n),\qquad \gamma_n-d_n\to\infty.
\end{equation}
Under~\eqref{eq:classical} and $\gamma_n=o(n)$, the oracle results also hold.
For $\gamma_n=(1+\epsilon)\log\log n$, ordinary variance consistency suffices.
\end{corollary}
\begin{proof}
Multiplying all costs by one positive scalar does not change their ordering.
Validity is equivalent to using the true $\sigma^2$ and the random effective
threshold $\gamma_n\widehat\sigma_n^2/\sigma^2$.
With probability tending to one this lies above
$d_n+(\gamma_n-d_n)/2$ and below $2\gamma_n$.
Use the lower deterministic threshold for true and oracle feasibility,
and the upper one for exclusion of partitions missing a boundary.
The proofs of Theorems~\ref{thm:main} and~\ref{thm:oracle} then apply unchanged.
\end{proof}

The variance corollary is conditional on properties of the estimate; it does
not assume that a global sample variance is consistent under changes.
No extension to dependent errors, arbitrary exponential families, or robust
rank tests is asserted without corresponding probability bounds.

\section{Multiscale validity with both subsegment tests}\label{sec:multiscale}

The plateau in Proposition~\ref{prop:plateau} suggests testing smaller intervals
inside a candidate segment. This avoids dilution by observations far from a
change, but the enlarged family of tests needs its own calibration. We retain
the Gaussian model~\eqref{eq:model}, the cost~\eqref{eq:cost}, and the
lexicographic objective. Only validity changes. In this section write
$\widehat\tau_n^{\mathrm{MS}}$ and $\widehat K_n^{\mathrm{MS}}$ for the new
estimator and its segment count.

\subsection{The effective length and the proposed threshold}

For positive integers $l,r$ with $l+r\le n$, define
\begin{equation}\label{eq:ms-scale}
h(l,r)=\frac{lr}{l+r},\qquad
p_n(l,r)=\sqrt{2\log\!\left(\frac{n(l+r)}{lr}\right)}.
\end{equation}
Here $h$ is the effective length of the two-sample contrast. It lies between
$\min(l,r)/2$ and $\min(l,r)$, and $h\le n/4$.
For $q\ge0$ and a candidate segment $(s,t]$, put
\begin{equation}\label{eq:ms-statistic}
T_n(Y;s,t)=\max_{s\le a<u<b\le t}
\bigl\{|Z_u(a,b)|-p_n(u-a,b-u)\bigr\}.
\end{equation}
Set the maximum to $-\infty$ for a singleton. A segment is multiscale valid
when $T_n(Y;s,t)\le q$, and $\widehat\tau_n^{\mathrm{MS}}$ minimizes $(K,Q)$
over partitions into such segments, with the previous tie convention.
In the units of~\eqref{eq:cost}, each split is therefore tested at
\begin{equation}\label{eq:ms-penalty}
\boxed{\displaystyle
\gamma_n(l,r;q)=\frac12\left[
\sqrt{2\log\!\left(\frac{n(l+r)}{lr}\right)}+q\right]^2.}
\end{equation}
This is a threshold for a split gain. The objective $Q$ receives no additional
penalty. Since the threshold depends on the split, the correction must be
applied before maximizing over $u$. Expanding the square gives
$\gamma_n=\log(n/h)+q\sqrt{2\log(n/h)}+q^2/2$; dropping the cross term changes
the calibration.

The size $n$ in~\eqref{eq:ms-scale} is the full, fixed sample size, even when
computing optimal partitions of prefixes. A balanced window with $l=r=r_0$
has correction $\sqrt{2\log(2n/r_0)}$. Thus long balanced windows are tested
at a lower threshold than short windows. A correction depending only on
$l+r$ would leave the endpoint splits of a long interval uncorrected:
on the full interval it would be constant, while the uncorrected maximum
absolute CUSUM diverges at the $\sqrt{\log\log n}$ scale.

\subsection{Simultaneous calibration and its sharp features}

For $n=1$ the singleton partition is always selected. All finite-sample
statements below concern $n\ge2$.
For $n\ge2$ define the pure-noise maximum and its upper tail by
\begin{equation}\label{eq:ms-null}
M_n=T_n(\sigma\varepsilon;0,n),\qquad
\alpha_n(q)=\Prob(M_n>q).
\end{equation}
The law does not depend on $\sigma$. The following is the only new imported
probability estimate.

\begin{lemma}[Simultaneous CUSUM control]\label{lem:ms-noise}
There are universal positive constants $c_0,c_1$ such that, for every
$n\ge2$ and $x>0$,
\begin{equation}\label{eq:ms-tail}
\Prob(M_n>c_0+c_1\sqrt{x})\le e^{-x}.
\end{equation}
In particular $M_n$ is uniformly bounded above in probability. For
$0<\alpha<1$, let
\begin{equation}\label{eq:ms-quantile}
q_{n,\alpha}=
\max\bigl\{0,\inf\{z\in\R:\Prob(M_n\le z)\ge1-\alpha\}\bigr\}.
\end{equation}
Then $\alpha_n(q_{n,\alpha})\le\alpha$ and
$q_{n,\alpha}\le c_0+c_1\sqrt{\log(1/\alpha)}$.
\end{lemma}
\begin{proof}
Apply \cite[Lemma 9, equation (79), p.~45 of the arXiv version]{verzelen2020}
to the independent standard Gaussian errors. Its interval convention
$[t_1,t_3)$ becomes ours under $(t_1,t_2,t_3)=(a+1,u+1,b+1)$.
Their contrast has the opposite sign, which disappears after squaring.
For $s=\log(n/h)\ge\log4$, that lemma bounds all squared noise contrasts,
with probability at least $1-e^{-x}$, by
\[
2s+C_1\log(s+x)+C_2x+C_3.
\]
The constants can be increased to be positive. Since
$\log(s+x)\le\log s+x/s$ and
$\sup_{s\ge\log4}(\log s)/\sqrt{s}<\infty$, taking square roots and
subtracting $\sqrt{2s}$ gives a bound $c_0+c_1\sqrt{x}$ uniformly in $s$.
This is also the calibration stated in their equation (40), p.~24.
The maximum in~\eqref{eq:ms-null} is over finitely many continuous Gaussian
absolute contrasts, so its distribution has no atoms. The quantile claims
follow, and clipping a negative quantile at zero can only reduce its tail.
\end{proof}

\begin{proposition}[Null consistency and the leading scale]\label{prop:ms-null}
Under a constant mean, for every deterministic $q_n\ge0$,
\begin{equation}\label{eq:ms-null-iff}
\Prob(\widehat K_n^{\mathrm{MS}}=1)\to1
\quad\Longleftrightarrow\quad q_n\to\infty.
\end{equation}
If instead all subinterval CUSUM gains are compared with the same scalar
$\gamma_n$, their pure-noise maximum divided by $\log n$ tends to one in
probability. Thus $\gamma_n=c\log n$ is null consistent for $c>1$ and
selects extra segments with probability tending to one for $c<1$.
Finally, replacing $p_n$ by $c p_n$, $0\le c<1$, in~\eqref{eq:ms-statistic}
fails under the null whenever $q_n=o(\sqrt{\log n})$.
\end{proposition}
\begin{proof}
The one-segment partition is selected exactly when $M_n\le q_n$.
Sufficiency follows from Lemma~\ref{lem:ms-noise}. For necessity, use the
single full-interval contrast at $u=\lfloor n/2\rfloor$. Its null law is
standard normal and its correction is at most $2$. Consequently
$\alpha_n(q_n)\ge2\Phi(-q_n-2)$, which stays positive on a subsequence if
$q_n$ does not tend to infinity.

For disjoint pairs of observations, the contrasts
$Z_{2i-1}(2i-2,2i)$, $1\le i\le\lfloor n/2\rfloor$, are independent
standard normals. Their maximum absolute value divided by
$\sqrt{2\log n}$ tends to one: Gaussian tail bounds and a union bound give
the upper bound, and the product of their acceptance probabilities gives
the lower bound. On the other hand, $h\ge1/2$ for every integer split,
so the maximum over all triples is at most
$\sqrt{2\log(2n)}+(M_n)_+$.
Lemma~\ref{lem:ms-noise} therefore gives the same limit for the full maximum.
Squaring and dividing by two proves the scalar-threshold claim.
For each disjoint pair, $p_n=\sqrt{2\log(2n)}$; the same independent maximum
exceeds $c p_n+q_n$ with probability tending to one when $c<1$.
\end{proof}

A fixed significance level gives controlled overestimation, whereas count
consistency requires an unbounded sequence $q_n$. For example, choosing
$\alpha_n^\circ=1/\log(en)$ and $q_n=q_{n,\alpha_n^\circ}$ ensures vanishing
null error and $q_n^2=O(\log\log(en))$. Slower divergence is also possible:
any deterministic $q_n\to\infty$ suffices for the null, provided the signal
conditions below hold. The constants in~\eqref{eq:ms-tail} are theoretical
universal constants; they are not proposed numerical tuning values.

For practical Gaussian calibration, simulate the maximum $M_n$ itself.
An ordinary empirical quantile is an approximation to~\eqref{eq:ms-quantile}.
Finite simulation error can be included explicitly: from $B$ independent
replicates, use the $k$th order statistic, clipped at zero, where
\[
\Prob\{\operatorname{Bin}(B,1-\alpha)\ge k\}\le\eta.
\]
Then, with calibration probability at least $1-\eta$, its true null tail
is at most $\alpha$. If no $k\le B$ satisfies this inequality, increase $B$.
With simulations independent of the observations, the unconditional
overestimation bound is at most $\alpha+\eta$. This recipe specifies
calibration; no simulated value is used in the proofs.

\subsection{Finite-sample recovery from local witnesses}

The local argument now uses a window that is itself among the tested
subintervals. It therefore has no factor depending on the size of the
candidate segment containing that window.

\begin{theorem}[Count and asymmetric localization]\label{thm:ms-finite}
Fix $n$, $q\ge0$, and the model~\eqref{eq:model}; abbreviate
$K=K_n^\star$ and suppress $n$ in the true lengths and jumps. Uniformly over
all such signals,
\begin{equation}\label{eq:ms-over}
\Prob(\widehat K_n^{\mathrm{MS}}>K)\le\alpha_n(q).
\end{equation}
For $K\ge2$, choose positive integer window lengths $L_j,R_j$ satisfying
\begin{equation}\label{eq:ms-windows}
L_j\le m_j,\quad R_j\le m_{j+1},\qquad
R_j+L_{j+1}\le m_{j+1}+1\quad(1\le j<K-1).
\end{equation}
Define the disjoint sets of possible boundaries
$\mathcal N_j=(\taus_j-L_j,\taus_j+R_j)\cap\Z$, and put
\begin{align}
h_j&=\frac{L_jR_j}{L_j+R_j},\qquad
A_j=\frac{\Delta_j\sqrt{h_j}}{\sigma},\qquad
B_j=\sqrt{2\log(n/h_j)}+q,\label{eq:ms-AB}\\
v_j&=\Phi(B_j-A_j)-\Phi(-B_j-A_j).
\label{eq:ms-v}
\end{align}
Then
\begin{align}
\Prob(\widehat K_n^{\mathrm{MS}}<K)&\le\sum_{j=1}^{K-1}v_j,
\label{eq:ms-under}\\
\Prob\left(\widehat K_n^{\mathrm{MS}}=K,\
\widehat\tau_{j,n}^{\mathrm{MS}}\in\mathcal N_j\ \text{for all }j<K\right)
&\ge1-\alpha_n(q)-\sum_{j=1}^{K-1}v_j.
\label{eq:ms-recovery}
\end{align}
Each $v_j\le\exp\{-(A_j-B_j)_+^2/2\}$. If the stronger inequalities
$R_j+L_{j+1}\le m_{j+1}$ hold, the sums of $v_j$ in these two bounds
can be replaced by $1-\prod_{j=1}^{K-1}(1-v_j)$.
\end{theorem}
\begin{proof}
On $\{M_n\le q\}$ every subinterval of every true segment passes its
test, since each such contrast has zero mean. The true partition is feasible;
minimality of the count gives~\eqref{eq:ms-over}.

If a partition contains no boundary in $\mathcal N_j$, the observations
$(\taus_j-L_j,\taus_j+R_j]$ lie in a single one of its segments.
Boundaries at the two window endpoints are allowed. Validity then requires
\[
|Z_{\taus_j}(\taus_j-L_j,\taus_j+R_j)|\le B_j.
\]
The two sides of this contrast have constant means and respective lengths
$L_j,R_j$. After a possible sign change it has law $\mathcal N(A_j,1)$,
so the probability of the displayed event is exactly $v_j$.
Outside the union of these events, every valid partition contains at least
one boundary in each $\mathcal N_j$. Condition~\eqref{eq:ms-windows} makes
these integer sets disjoint, hence the count is at least $K$. This proves
\eqref{eq:ms-under}. Intersecting with $\{M_n\le q\}$ gives equality of
counts and exactly one boundary in each ordered neighborhood, proving
\eqref{eq:ms-recovery}. No independence from the true-feasibility event
is needed.

If $A_j>B_j$, the acceptance event is contained in a standard Gaussian
lower-tail event at $B_j-A_j$; its probability is at most
$\exp\{-(A_j-B_j)^2/2\}$. For $A_j\le B_j$ use the bound one.
Under the stronger spacing inequalities the observation windows are
disjoint, so their acceptance events are independent and their union
has probability $1-\prod_j(1-v_j)$.
\end{proof}

\begin{corollary}[A radius with prescribed confidence]\label{cor:ms-radius}
Let $K\ge2$, $q=q_{n,\alpha}$, and $0<\beta<1$. Choose integer radii $r_j$
with $1\le r_j\le\min(m_j,m_{j+1})$ and
$r_j+r_{j+1}\le m_{j+1}+1$. If
\begin{equation}\label{eq:ms-radius}
\frac{\Delta_j}{\sigma}\sqrt{\frac{r_j}{2}}
\ge \sqrt{2\log(2n/r_j)}+q+
\sqrt{2\log((K-1)/\beta)}\qquad(j<K),
\end{equation}
then the count is correct and every ordered location error is less than
$r_j$, simultaneously, with probability at least $1-\alpha-\beta$.
\end{corollary}
\begin{proof}
Use $L_j=R_j=r_j$ in Theorem~\ref{thm:ms-finite}. Each tail bound is at
most $\beta/(K-1)$.
\end{proof}

The smallest admissible radius can be found from the monotone inequality
\eqref{eq:ms-radius}. A simpler sufficient inequality is
\begin{equation}\label{eq:ms-radius-sufficient}
\frac{r_j\Delta_j^2}{\sigma^2}
\ge12\log(2n/r_j)+6q^2+12\log((K-1)/\beta),
\end{equation}
obtained by $(a+b+c)^2\le3(a^2+b^2+c^2)$.
This yields the scale $\sigma^2/\Delta_j^2$ multiplied by local logarithmic
and confidence terms, subject to the window spacing conditions. In
comparison, the plain-SVP validity bound in Corollary~\ref{cor:radii}
contains a term proportional to $\sigma\sqrt{n\gamma_n}/\Delta_n$.
Both are guarantees obtained from validity alone; neither is the final
fixed-signal oracle localization rate.

\subsection{Consistency with short segments and shrinking jumps}

\begin{theorem}[Local signal consistency]\label{thm:ms-consistency}
Let $q_n\ge0$ be deterministic with $q_n\to\infty$. Suppose $K_n^\star\ge2$,
$\ell_n\to\infty$, and, with $d_{j,n}=\min(m_{j,n},m_{j+1,n})$,
\begin{equation}\label{eq:ms-local-snr}
\min_{j<K_n^\star}
\frac{d_{j,n}\Delta_{j,n}^2/\sigma^2}
{\log(en/d_{j,n})+q_n^2+\log(K_n^\star+1)}\longrightarrow\infty.
\end{equation}
Then, for every $\eta>0$,
\begin{equation}\label{eq:ms-relative}
\Prob\left(\widehat K_n^{\mathrm{MS}}=K_n^\star,\quad
\max_{j<K_n^\star}
\frac{|\widehat\tau_{j,n}^{\mathrm{MS}}-\taus_j|}{d_{j,n}}<\eta\right)
\longrightarrow1.
\end{equation}
For fixed $K^\star\ge2$ and $\ell_n\ge an$, it is sufficient that
\begin{equation}\label{eq:ms-macro}
q_n\to\infty,\qquad q_n^2=o(n\Delta_n^2/\sigma^2).
\end{equation}
For $K_n^\star=1$, the necessary and sufficient condition is simply
$q_n\to\infty$.
\end{theorem}
\begin{proof}
It suffices to take $0<\eta<1/2$. Choose
$r_{j,n}=\lfloor\eta d_{j,n}\rfloor$. For all large $n$ these are positive
and at least $\eta d_{j,n}/2$. Since both adjacent $d_{j,n}$ are at most
their common segment length, their radii satisfy~\eqref{eq:ms-windows}.
Let $D_{j,n}$ be the denominator in~\eqref{eq:ms-local-snr}, and let
$E_n$ be the minimum ratio in that display. For the quantities in
Theorem~\ref{thm:ms-finite}, there is a constant $C_\eta$ such that
\[
A_{j,n}^2\ge\frac{\eta}{4}E_nD_{j,n},\qquad
B_{j,n}^2\le C_\eta D_{j,n}.
\]
The second inequality uses
$\log(2n/r_{j,n})\le\log(en/d_{j,n})+\log(4/(e\eta))$.
Since $E_n\to\infty$, eventually $B_{j,n}\le A_{j,n}/2$ uniformly in $j$.
Thus
\[
\sum_{j<K_n^\star}v_{j,n}
\le K_n^\star\exp\!\left\{-\frac{\eta E_n}{32}
\bigl(q_n^2+\log(K_n^\star+1)\bigr)\right\}\longrightarrow0.
\]
Lemma~\ref{lem:ms-noise} gives $\alpha_n(q_n)\to0$.
Theorem~\ref{thm:ms-finite} now proves~\eqref{eq:ms-relative}.
In the fixed-count macroscopic regime, $D_{j,n}=O(1+q_n^2)$ and
$d_{j,n}\ge an$, giving~\eqref{eq:ms-macro}. The no-change assertion is
Proposition~\ref{prop:ms-null}.
\end{proof}

The condition $\ell_n\to\infty$ is used for the relative-error conclusion.
Finite-sample Theorem~\ref{thm:ms-finite} also covers length-one segments:
choosing $L_j=R_j=1$ gives $\mathcal N_j=\{\taus_j\}$ and an exact-recovery
bound when the jumps are sufficiently large. The $+1$ in the integer
spacing condition permits consecutive true boundaries in this case.

For macroscopic segments, the correction at a local balanced window is
bounded. There is consequently no compulsory $\log\log n$ factor in
\eqref{eq:ms-macro}. The sequence $q_n$ must nevertheless be specified:
no one prescribed diverging sequence adapts to every arbitrarily slowly
diverging signal energy. For a fixed sequence, its squared size must remain
small compared with the signal energy. Under fixed error probabilities,
the finite-sample bounds apply with their stated nonvanishing error level.

\subsection{Oracle equivalence and the meaning of faster convergence}

\begin{theorem}[Multiscale oracle equivalence]\label{thm:ms-oracle}
Under the fixed-signal regime~\eqref{eq:classical}, if
$q_n\to\infty$ and $q_n^2=o(n)$, then
\begin{equation}\label{eq:ms-oracle}
\Prob(\widehat\tau_n^{\mathrm{MS}}=\tol)\longrightarrow1.
\end{equation}
The joint location, level and empirical-loss limits in
Theorem~\ref{thm:limits} hold for multiscale SVP with the same limiting laws.
\end{theorem}
\begin{proof}
Count consistency follows from Theorem~\ref{thm:ms-consistency} (or
Proposition~\ref{prop:ms-null} when $K^\star=1$). We prove oracle feasibility
separately. Fix $M\ge0$ and consider all intervals whose endpoints are
within $M$ of the endpoints of a true segment. After subtracting that
segment's level, each interval has a deterministic contamination vector
supported on at most $2M$ observations. Under~\eqref{eq:classical} its
Euclidean norm divided by $\sigma$ is at most a constant $C_M$.
Every standardized contrast has coefficient vector of Euclidean norm one;
Cauchy--Schwarz therefore gives, simultaneously for all these intervals,
\[
T_n(Y;a,b)\le M_n+C_M.
\]
Lemma~\ref{lem:ms-noise} and $q_n\to\infty$ make the right-hand side at
most $q_n$ with probability tending to one for fixed $M$.
Lemma~\ref{lem:oracle-tight} places the oracle endpoints in this family with
probability arbitrarily close to one by taking $M$ large. Thus the oracle
is multiscale feasible with probability tending to one. At the correct
count, minimizing $Q$ over the valid subset containing its unrestricted
minimizer gives~\eqref{eq:ms-oracle}. The oracle distributional results
then transfer exactly as in Theorem~\ref{thm:limits}.
\end{proof}

This feasibility proof only needs a bounded contamination argument on the
absolute-contrast scale. The endpoint-layer analysis needed for the
second-order plain threshold disappears. The improvement is therefore
both statistical and in the proof: weak or short changes satisfy less
restrictive signal conditions, and the simultaneous validity argument is
shorter. For fixed nonzero jumps on macroscopic segments, both estimators
already attain $O_{\Prob}(1)$ error in observation indices, or
$O_{\Prob}(n^{-1})$ after normalization. The information bound following
Theorem~\ref{thm:limits} precludes a uniformly vanishing integer error at
fixed jump size. We make no claim of a faster rate in that regime, nor of
oracle equivalence throughout the shrinking-jump regime of
Theorem~\ref{thm:ms-consistency}.

\begin{corollary}[Recovery of the short plateau]\label{cor:ms-plateau}
For the plateau of Proposition~\ref{prop:plateau}, suppose
$m_n\to\infty$, $m_n=o(n)$, $q_n\to\infty$, and
\begin{equation}\label{eq:ms-plateau}
\log(en/m_n)+q_n^2=o(m_n\Delta^2/\sigma^2).
\end{equation}
Then $\widehat K_n^{\mathrm{MS}}=3$ with probability tending to one and
both ordered location errors are $o_{\Prob}(m_n)$.
In particular, the fixed-height plateau with $m_n\asymp n^{1/4}$ missed by
plain SVP is recovered by multiscale SVP if $q_n^2=o(n^{1/4})$.
\end{corollary}
\begin{proof}
For large $n$ both adjacent minimum lengths equal $m_n$ and $K^\star=3$.
Apply Theorem~\ref{thm:ms-consistency}.
\end{proof}

The contrast on a balanced window of radius $r\le m_n$ around a plateau
edge has squared mean $r\Delta^2/(2\sigma^2)$. Its signal therefore grows
linearly in the window length, while the whole-series scan in
\eqref{eq:plateau} has signal of order $m_n^2\Delta^2/(\sigma^2n)$.
This calculation identifies the source of the gain without relying on a
simulation or on a claim that the finite-sample constants are minimax optimal.

\subsection{Hereditary validity, implementation and the relation to SMUCE}

\begin{proposition}[Stability for both pruning rules]\label{prop:ms-pruning}
For fixed $n$ and $q$, multiscale invalidity propagates from an interval to
every interval containing it. The lexicographic dynamic program with this
validity constraint therefore permits both the left and right validity
pruning rules.
\end{proposition}
\begin{proof}
If $[s,t]\subseteq[s',t']$, every triple tested in~\eqref{eq:ms-statistic}
for $(s,t]$ is also tested for $(s',t']$, with exactly the same threshold.
Thus $T_n(Y;s,t)\le T_n(Y;s',t')$. If an interval starting at boundary $s$
fails at endpoint $t$, every later endpoint with the same start fails.
Every older start at the same or any later endpoint also fails. Removing
these transitions cannot remove a feasible partition. The additive cost
and lexicographic optimal-substructure argument then give the exact
optimizer over multiscale-valid partitions.
\end{proof}

Equivalently, use the base interval test
\[
g_n(Y_{a..b})=\max_{a<u<b}
\{|Z_u(a,b)|-p_n(u-a,b-u)\}
\]
and enforce both subsegment rules. In a chronological implementation that
introduces every possible start, examines surviving starts from newest to
oldest, and deletes a failing start and all older starts, this realizes
\eqref{eq:ms-statistic} without explicitly storing all subinterval maxima.
Indeed, a failed subinterval is examined before any surviving interval
that contains it. If its start was already deleted, an earlier failed
subinterval has already invalidated every containing start. Induction over
endpoints proves that the retained starts are precisely those whose current
segments satisfy the aggregated constraint. This argument concerns the
outer validity pruning; it requires exact evaluation of each base test.

The present Gaussian implementation compares an uncorrected maximum split
gain with a scalar \texttt{gamma}. Simply selecting
\texttt{subtests = "both"} implements the scalar-threshold variant, whose
null scale is described in Proposition~\ref{prop:ms-null}.
To implement~\eqref{eq:ms-penalty}, the base validity test must account for
both split lengths. Existing functional pruning inside a CUSUM routine
must be justified again for this modified statistic. For an unstandardized
RSS gain, the threshold is $2\sigma^2\gamma_n(l,r;q)$; for a half-RSS gain
it is $\sigma^2\gamma_n(l,r;q)$.

The global $n$ must remain fixed during a run. Replacing it by the length
of the containing candidate segment changes the threshold of a fixed
witness when that segment grows, so the preceding nesting proof would
no longer apply. More effective outer pruning can reduce computation, but
the new split-dependent base test also has a cost. Statistical consistency
does not itself imply a runtime improvement.

Finally, raw maximization of the Gaussian single-change GLR over
subintervals is not the SMUCE statistic. In the Gaussian case, SMUCE
constrains, for a common proposed segment level $\theta$, quantities of the
form
\begin{equation}\label{eq:smuce-distinction}
\frac{\sqrt{|I|}}{\sigma}|\overline Y_I-\theta|
-\sqrt{2\log(en/|I|)}
\end{equation}
over all intervals $I$ inside that segment, and uses likelihood optimized
over the admissible levels~\cite{frick2014}. Our constraint compares two
adjacent empirical means and uses the unconstrained segment-mean cost
\eqref{eq:cost}. Both belong to multiscale minimum-count segmentation, but
their feasible sets and secondary costs must be distinguished. An exact
SVP--SMUCE identification requires SMUCE's own validity and constrained cost.

\begin{table}[ht]
\caption{Calibration and conclusions in the known-variance Gaussian model.
The null scales in the first two rows are leading orders.}\label{tab:ms-comparison}
\centering\small
\begin{tabular}{@{}>{\raggedright\arraybackslash}p{0.20\linewidth}p{0.27\linewidth}p{0.46\linewidth}@{}}
\toprule
Validity rule & Null calibration & Main statistical conclusion\\
\midrule
Plain scan & $\gamma_n\sim\log\log n$, with a diverging margin above $d_n$
& Macroscopic consistency; possible cancellation of short plateaus.\\[0.5em]
Both, scalar threshold & $\gamma_n=c\log n$, $c>1$
& Local windows are tested, but every scale pays the logarithmic threshold.\\[0.5em]
Both, multiscale & $\gamma_n(l,r;q_n)$ from~\eqref{eq:ms-penalty}
& Uniform overestimation control and local recovery bounds; $q_n\to\infty$
for count consistency.\\
\bottomrule
\end{tabular}
\end{table}

\appendix
\section{Oracle localization and the local experiment}\label{app:oracle}

For $K^\star=1$ the oracle partition has no estimated boundary, so localization
is immediate. In the rest of this appendix suppose $K^\star\ge2$.
We give details so that the oracle conclusions do not rest on an unexplained
appeal to classical change-point asymptotics. The single imported localization
input is the following special case of \cite[Theorem 3, pp.~1926--1927]{wang2020}:
under~\eqref{eq:classical}, a sufficiently large residual-sum-of-squares penalty
$C\sigma^2\log n$ yields the correct count and maximum boundary error at most
$C'\log n$ with probability tending to one. Their Assumption 2 holds because
$\Delta\sqrt{\ell_n}/\sigma\asymp\sqrt n$; their independent sub-Gaussian
noise assumption includes~\eqref{eq:model}. Their count denotes changes;
ours denotes segments. Their cost has no factor $1/(2\sigma^2)$.

On the event that this penalized estimator selects $K^\star$ segments, it
also minimizes $Q$ among all $K^\star$-segment partitions. Almost-sure
uniqueness therefore implies equality with $\tol$. We obtain a deterministic
$C_0<\infty$ such that
\begin{equation}\label{eq:oracle-log}
\Prob\left(\max_j|\widetilde\tau_j-\taus_j|\le C_0\log n\right)\to1.
\end{equation}
This use of penalized least squares is independent of SVP consistency.

\subsection{Uniform profile-likelihood expansion}
Let $h=(h_1,\ldots,h_{K^\star-1})\in\Z^{K^\star-1}$ with
$\norm h_\infty\le b_n=\lceil C_0\log n\rceil$, and set $h_0=h_{K^\star}=0$.
For large $n$, all resulting partitions $\taus+h$ are ordered and their
boundary neighborhoods are disjoint.
Let $I_k(h)=(\taus_{k-1}+h_{k-1},\taus_k+h_k]$ and define
\[
A_k(h)=\sum_{i\in I_k(h)}(Y_i-\theta_k),\qquad L_k(h)=|I_k(h)|,
\qquad D_n(h)=\frac{1}{2\sigma^2}\sum_k\frac{A_k(h)^2}{L_k(h)}.
\]
Completing squares on each interval gives
\begin{equation}\label{eq:profile}
Q(\taus+h)-Q(\taus)
=\sum_{j=1}^{K^\star-1}R_{j,n}(h_j)-\{D_n(h)-D_n(0)\},
\end{equation}
where the local walks are exactly
\[
R_{j,n}(h)=
\begin{cases}
\displaystyle\sum_{r=1}^{h}\left(\delta_j^2/2+
\delta_j\varepsilon_{\taus_j+r}\right),&h>0,\\
0,&h=0,\\
\displaystyle\sum_{r=1}^{-h}\left(\delta_j^2/2-
\delta_j\varepsilon_{\taus_j-r+1}\right),&h<0.
\end{cases}
\]
For example, moving boundary $j$ to the right assigns observations of level
$\theta_{j+1}$ to $\theta_j$; the per-observation cost difference in the
units of~\eqref{eq:cost} is $\delta_j^2/2+\delta_j\varepsilon_i$.
This verifies the signs and the factor of two.

We have $A_k(0)=\Op(\sqrt n)$, uniformly over the fixed number of segments.
By the maximal inequality for Gaussian partial sums over the finitely many
boundary windows,
\[
\max_{k,\norm h_\infty\le b_n}|A_k(h)-A_k(0)|
=\Op(b_n+\sqrt{b_n})=\Op(\log n),
\quad
\max_{k,h}|L_k(h)-L_k(0)|\le2b_n.
\]
The deterministic part of the first bound is the contribution of
misassigned fixed levels. All lengths are at least a positive multiple of $n$.
Expanding the squares and reciprocals therefore gives
\begin{equation}\label{eq:uniform-remainder}
\sup_{\norm h_\infty\le b_n}|D_n(h)-D_n(0)|
=\Op\left(\frac{b_n}{\sqrt n}+\frac{b_n^2}{n}+\frac{b_n}{n}\right)=\op(1).
\end{equation}

\subsection{Argmin tightness and distribution}
Within this box the walks $R_{j,n}$ are independent over $j$ and the two
sides, and their joint law is exactly the restriction of the independent
walks in~\eqref{eq:random-walk} to $[-b_n,b_n]$.
For these infinite walks the strong law gives $R_j(h)\to+\infty$ as
$h\to\pm\infty$. All differences between distinct walk values have
nondegenerate continuous distributions. Hence $H_j$ exists and is unique
almost surely. For the sum $\mathcal R(h)=\sum_jR_j(h_j)$ on
$\Z^{K^\star-1}$, there is also a unique minimum at $H=(H_j)_j$ and a
strictly positive random gap to the second smallest value: every bounded
sublevel set is finite, because each summand is bounded below and tends
to $+\infty$ at both ends.

As $b\to\infty$, the minimizer and the gap on $[-b,b]^{K^\star-1}$ agree
eventually with the unrestricted ones, almost surely.
For any fixed $\eta>0$, the probability that the finite-box gap is at most
$\eta$ thus approaches the corresponding probability for the infinite gap;
the latter tends to zero as $\eta\downarrow0$.
The uniform perturbation~\eqref{eq:uniform-remainder} therefore changes the
minimizer with probability tending to zero. Combining with
\eqref{eq:oracle-log} gives
\begin{equation}\label{eq:oracle-walk-limit}
(\widetilde\tau_j-\taus_j)_{j<K^\star}\Longrightarrow(H_j)_{j<K^\star},
\end{equation}
and proves Lemma~\ref{lem:oracle-tight}. This gap argument requires neither
independence of the profile remainder nor a continuous-parameter argmin theorem.

\subsection{Independence from level estimates}
Fix $M$. The walks restricted to $[-M,M]$ involve finitely many observations
near the true boundaries. Delete these observations from the normalized
true-segment sums. The remaining sums are independent of the local walks
and converge jointly to independent standard normals; deletion changes each
normalized sum by $\op(1)$. Thus the truncated walk minimizers and the
level sums have independent limits. The probability that any full walk
minimizer lies outside $[-M,M]$ tends to zero as $M\to\infty$.
The perturbation and localization arguments above transfer this statement
to the oracle estimator. Taking $n\to\infty$ and then $M\to\infty$ proves
the joint independence used in Theorem~\ref{thm:limits}.

\section{Comparison with the earlier consistency note}\label{app:comparison}

The earlier companion note establishes count consistency for fixed jumps
and $\ell_n\asymp n$ at $\gamma_n=c\log n$, $c>1$.
The principal differences are as follows.
\begin{enumerate}
\item The Gaussian dependence structure supplies the $\log\log n$ null
scale and its necessary second-order correction. Treating the splits as
unrelated tests obscures this scale.
\item Local windows around the true changes give a single finite-sample
event excluding every missing-boundary partition. There is no factor $n^3$
from a union over candidate triples, and the jump signs cause no ambiguity.
\item Count, localization, oracle equality, level inference, and empirical
loss have separate statements with their respective assumptions.
\item OP equivalence is a consequence of correct counts and common cost
units. The inequality is $K^{\mathrm{SVP}}\le K^{\mathrm{OP}}$; the earlier
note's introductory reference to SVP using strictly more segments has the
direction reversed. A sufficiently large logarithmic penalty is not a
theorem for every named BIC convention.
\item The central plateau explicitly rules out an unjustified extension
from macroscopic segments to arbitrary local spacing--jump regimes.
\item The multiscale extension in Section~\ref{sec:multiscale} repairs this
failure by changing validity. It imports a simultaneous CUSUM inequality
and derives explicit SVP recovery bounds, a local signal consistency theorem,
and a simpler fixed-signal oracle-feasibility proof. Its tuning parameter
$q$ is on the absolute-contrast scale and is different from the ordinary
GLR threshold $\gamma$.
\end{enumerate}
The finite-sample constants in the local-window bounds are explicit but
not claimed optimal. Sharpness is established for the null threshold and
the fixed-signal location scale; it is not a claim of uniform minimax
optimality of plain SVP. The proofs are mathematical arguments, not
machine-checked formalizations, and no priority claim is made.

\begin{thebibliography}{9}
\bibitem{torgovitski2016}
L.~Torgovitski.
\newblock A Darling--Erd\H{o}s-type CUSUM-procedure for functional data II.
\newblock arXiv:1407.3625v2, 2016.
\newblock \url{https://arxiv.org/abs/1407.3625}.
\bibitem{wang2020}
D.~Wang, Y.~Yu and A.~Rinaldo.
\newblock Univariate mean change point detection: penalization, CUSUM and optimality.
\newblock \emph{Electronic Journal of Statistics}, 14(1):1917--1961, 2020.
\newblock \url{https://doi.org/10.1214/20-EJS1710}.
\bibitem{verzelen2020}
N.~Verzelen, M.~Fromont, M.~Lerasle and P.~Reynaud-Bouret.
\newblock Optimal change-point detection and localization.
\newblock \emph{Annals of Statistics}, 51(4):1586--1610, 2023.
\newblock \url{https://doi.org/10.1214/23-AOS2297}.
\newblock Detailed statements and proofs cited here use arXiv:2010.11470v2,
2020, \url{https://arxiv.org/abs/2010.11470v2}.
\bibitem{frick2014}
K.~Frick, A.~Munk and H.~Sieling.
\newblock Multiscale change point inference.
\newblock \emph{Journal of the Royal Statistical Society: Series B},
76(3):495--580, 2014.
\newblock \url{https://doi.org/10.1111/rssb.12047}.
\end{thebibliography}
\end{document}
